\subsection*{2. A ``jeu de Légo-Teichmüller'' and the Galois group}
The second section develops the program from the shared geometric
intuition of surfaces that Grothendieck describes as emerging from his
teaching (pp.~2--3). The transition from combinatorial topology to complex
algebraic curves brings back the fields, rings, and groups named on p.~3,
then the ``multiplicités modulaires'' and their automorphisms. On
pp.~4--5 the family of moduli multiplicities, its gluing and erasing
operations, and the corresponding fundamental groupoids form the
``tour de Teichmüller'', on which he considers the action of the absolute
Galois group. He states the faithfulness of the action on the first
nontrivial level and asks for a description of the tower's action in terms
of that level; p.~5 explicitly says that he has no conjecture yet for a
characterization of the Galois group inside the outer automorphism group.
Note (1), added on p.~52, later withdraws the accompanying judgment
``hors de portée'' and reports a renewed step toward that characterization.
The groupoid language keeps the chosen collections of base points visible
under the symmetries. On pp.~6--7 he formulates the construction principle
in which the first two moduli dimensions supply generators and relations,
sets it beside Demazure's rank principle, and describes the remaining
topological work in terms of the fundamental group at infinity and
``dévissage'' of stratified structures, with an explicit reference to
\S5. On pp.~8--9 the ``petite boite'' and ``grande boite'' turn those
operations into assemblies of surfaces; Thurston's surgery is named,
canonical cellular decompositions point forward to \S3, and real
structures guide the proposed decompositions of moduli multiplicities.
A mathematical slip remains printed on p.~9: the relations
$\rho^3=\sigma^2$, $\rho^4=\sigma^6=1$ imply that
$z=\rho^3=\sigma^2$ satisfies $z^4=z^3=1$, hence $z=1$,
$\rho=1$, and $\sigma^2=1$. The resulting presentation has mutually
inverse maps to the cyclic group of order two, sending $\rho$ to its
identity and $\sigma$ to its generator, and sending that generator back
to $\sigma$. In contrast, the map
$n\longmapsto M_n=\left(\begin{smallmatrix}1&n\\0&1\end{smallmatrix}\right)$
from $\mathbb Z$ to $\mathrm{SL}_2(\mathbb Z)$ is an injective
homomorphism: $M_mM_n=M_{m+n}$, and its upper-right entry recovers $n$.
This establishes the
failure of the displayed presentation to present the group named there.
The readers retain the scanned exponents. For subsequent work on the
Galois-group program, one may name \emph{La Longue Marche à travers la
théorie de Galois} and Drinfeld's
\href{https://www.mathnet.ru/eng/aa199}{\emph{On quasitriangular quasi-Hopf algebras and on a group that is closely connected with}
$\mathrm{Gal}(\overline{\mathbb Q}/\mathbb Q)$}.
(\href{https://grothendieck.umontpellier.fr/web/viewer.html?file=../119.pdf\#page=5}{Copy 1, pp.~2--9, beginning at site 4};
\href{https://grothendieck.umontpellier.fr/web/viewer.html?file=../119.pdf\#page=12}{mathematical slip on p.~9, site 11}.)

